% Chapter: Proof Theory Automation
% Part of the CCTT Monograph

\chapter{Proof Theory: Automation and Tactics}
\label{ch:proof-theory-automation}

While the proof theory of the previous chapters provides a language for
\emph{writing} proofs, humans (and LLMs) benefit from automation.
This chapter describes the tactic language: functions that
\emph{construct} proof terms automatically.

The key design principle: \textbf{tactics are optional}.  A human can
always write a proof term directly.  Tactics just make it easier by
automating routine steps.  Every tactic outputs a proof term that
the checker can independently verify.

\section{Tactic Architecture}

A tactic is a function
\[
  \tau : \mathsf{OTerm} \times \mathsf{OTerm} \times \mathsf{Context}
  \to \mathsf{TacticResult}
\]
where $\mathsf{TacticResult}$ contains either a proof term (success)
or a failure reason.

Tactics are classified by their power and cost:

\begin{enumerate}
\item \textbf{Finishing tactics}: close the goal immediately.
  \begin{itemize}
  \item $\mathsf{refl}$: try reflexivity (syntactic equality up to
    normalization).
  \item $\mathsf{assumption}$: match an in-scope hypothesis.
  \item $\mathsf{definitional}$: normalize both sides and compare.
  \end{itemize}

\item \textbf{Simplification tactics}: transform the goal.
  \begin{itemize}
  \item $\mathsf{simp}(\mathcal{A})$: apply simplification axioms from
    set $\mathcal{A}$ exhaustively.
  \item $\mathsf{ring}$: prove ring equalities by polynomial
    normalization.
  \item $\mathsf{omega}$: prove linear arithmetic via Z3 QF\_LIA.
  \end{itemize}

\item \textbf{Decomposition tactics}: break the goal into subgoals.
  \begin{itemize}
  \item $\mathsf{induction}(x, \mathit{structure})$: set up induction on
    $x$ over naturals or lists.
  \item $\mathsf{cases}(d, \mathit{labels})$: case-split on discriminant $d$.
  \item $\mathsf{ext}(x)$: prove function equality by extensionality.
  \end{itemize}

\item \textbf{Rewriting tactics}: apply equations.
  \begin{itemize}
  \item $\mathsf{unfold}(f)$: $\delta$-expand a function definition.
  \item $\mathsf{rewrite}(e, \mathit{dir})$: apply a proven equation
    as a rewrite.
  \end{itemize}

\item \textbf{Strategy tactics}: set up complex proof structures.
  \begin{itemize}
  \item $\mathsf{by\_simulation}(R)$: construct a simulation proof skeleton.
  \item $\mathsf{by\_spec}(S)$: construct an abstraction refinement.
  \item $\mathsf{by\_loop\_invariant}(I)$: construct a loop invariant proof.
  \item $\mathsf{by\_comm\_diagram}$: construct a commutative diagram proof.
  \end{itemize}

\item \textbf{Search tactics}: try multiple approaches.
  \begin{itemize}
  \item $\mathsf{auto}(\mathit{depth})$: automatic proof search up to a
    given depth.
  \end{itemize}
\end{enumerate}


\section{The \texorpdfstring{$\mathsf{auto}$}{auto} Tactic}

The $\mathsf{auto}$ tactic is a depth-bounded proof search that tries
tactics in order of increasing complexity:

\begin{enumerate}
\item Reflexivity (identity)
\item Definitional equality (normalization)
\item Arithmetic simplification (constant folding, Z3)
\item Assumption matching
\item Congruence (recursive on arguments)
\item Symmetry + recursive search
\end{enumerate}

\begin{lstlisting}
def auto(lhs, rhs, ctx, depth=3):
    # 1. Reflexivity
    if terms_equal(normalize(lhs), normalize(rhs)):
        return Refl(lhs)
    # 2. Arithmetic
    if is_constant(lhs) and is_constant(rhs):
        if eval(lhs) == eval(rhs):
            return ArithmeticSimp(lhs, rhs)
    # 3. Congruence (recursive)
    if head(lhs) == head(rhs) and depth > 0:
        proofs = [auto(a, b, ctx, depth-1)
                  for a, b in zip(args(lhs), args(rhs))]
        if all_success(proofs):
            return Cong(head(lhs), proofs)
    # 4. Symmetry
    if depth >= 2:
        p = auto(rhs, lhs, ctx, depth-1)
        if success(p): return Sym(p)
    # Failure
    return None
\end{lstlisting}

\begin{remark}
The $\mathsf{auto}$ tactic has worst-case exponential time in
$\mathit{depth}$, but with depth $\leq 3$ it is fast in practice.
For deep proofs, the user should guide the search with explicit
proof terms or intermediate lemmas.
\end{remark}


\section{The \texorpdfstring{$\mathsf{ring}$}{ring} Tactic}

The $\mathsf{ring}$ tactic proves equalities in commutative rings
by normalizing both sides to multivariate polynomial form.

\textbf{Algorithm}:
\begin{enumerate}
\item Convert both sides to polynomials: sums of products of variables
  with integer coefficients.
\item A polynomial is a map from monomials to coefficients.
\item A monomial is a sorted tuple of (variable, exponent) pairs.
\item Two polynomials are equal iff they have the same monomials with
  the same coefficients.
\end{enumerate}

\begin{example}
$\mathsf{ring}$ can prove:
\[
  (a + b)^2 = a^2 + 2ab + b^2
\]
by expanding both sides to the polynomial $\{a^2: 1, ab: 2, b^2: 1\}$.
\end{example}

If polynomial normalization fails (e.g., division or non-polynomial
operations), $\mathsf{ring}$ falls back to Z3 QF\_NIA.


\section{The \texorpdfstring{$\mathsf{omega}$}{omega} Tactic}

The $\mathsf{omega}$ tactic handles linear integer arithmetic:
\[
  \sum_i a_i x_i + c \;\bowtie\; 0
  \quad \text{where } \bowtie \;\in\; \{=, \neq, <, \leq, >, \geq\}
\]

Implementation: construct a Z3 $\mathsf{QF\_LIA}$ formula and check
validity.  If Z3 returns UNSAT for the negation, the formula is valid
and we produce a $\mathsf{Z3Discharge}$ proof term.

\begin{remark}
Z3 is used here as a \emph{decision procedure}, not as a sampler.
A Z3 UNSAT result is a machine-checked proof.  If Z3 times out,
$\mathsf{omega}$ reports failure---it never returns ``probably true.''
\end{remark}


\section{The \texorpdfstring{$\mathsf{simp}$}{simp} Tactic}

The simplification tactic applies a set of rewrite rules exhaustively.
Given a set of axioms $\mathcal{A} = \{\ell_i \equiv r_i\}$:

\begin{enumerate}
\item Normalize both sides ($\beta, \delta$ reduction + constant folding).
\item For each axiom $\ell_i \equiv r_i$, try to match $\ell_i$ against
  subterms of the goal.  If a match is found, rewrite.
\item Repeat until no more rewrites apply.
\item Check if normalized sides are equal.
\end{enumerate}

The output is a $\mathsf{RewriteChain}$ proof term recording each
rewrite step.

\begin{remark}
Termination of $\mathsf{simp}$ is ensured by a maximum iteration
count (100).  Confluence is not guaranteed for arbitrary axiom sets,
but the CCTT axioms are designed to be terminating and confluent
on normal forms.
\end{remark}


\section{Tactic Combinators}

Tactics compose via combinators:

\begin{definition}[Sequencing]
$\mathsf{then}(\tau_1, \tau_2)$: try $\tau_1$; if it fails, try $\tau_2$.
\end{definition}

\begin{definition}[Alternative]
$\mathsf{orelse}(\tau_1, \ldots, \tau_n)$: try each tactic in order;
return the first success.
\end{definition}

\begin{definition}[Repetition]
$\mathsf{repeat}(\tau, \mathit{max})$: apply $\tau$ repeatedly (up to
$\mathit{max}$ times) until it fails or the goal is solved.
\end{definition}

These combinators always produce valid proof terms because each
constituent tactic produces a valid term, and the combinators compose
them via $\mathsf{Trans}$ or $\mathsf{RewriteChain}$.


\section{Integration with Z3}

Z3 appears in the proof theory in two distinct roles:

\begin{enumerate}
\item \textbf{Decision procedure} (via $\mathsf{Z3Discharge}$):
  For quantifier-free formulas in decidable fragments, Z3 provides
  a definitive YES/NO answer.  This is a \emph{proof}, not sampling.

\item \textbf{Tactic backend} (via $\mathsf{omega}$ and $\mathsf{ring}$):
  Tactics use Z3 to discharge arithmetic subgoals.  The resulting
  proof term records the Z3 call; re-verification repeats it.
\end{enumerate}

In both cases, Z3 timeout means proof \textbf{rejection}.  This is
the critical design decision that distinguishes our approach from
sampling-based methods.


\section{What Can and Cannot Be Automated}

\begin{longtable}{p{6cm}cc}
\toprule
\textbf{Task} & \textbf{Automated?} & \textbf{Tactic} \\
\midrule
Constant arithmetic & Yes & $\mathsf{omega}$ \\
Polynomial equalities & Yes & $\mathsf{ring}$ \\
Syntactically identical programs & Yes & $\mathsf{auto}$ \\
Programs differing by commutativity & Yes & $\mathsf{simp}$ \\
Same algorithm, different variable names & Yes & $\mathsf{refl}$ (after normalization) \\
\midrule
Recursive vs.\ iterative & Skeleton & $\mathsf{induction}$ \\
Different algorithms (same spec) & Skeleton & $\mathsf{by\_spec}$ \\
State-space correspondence & Skeleton & $\mathsf{by\_simulation}$ \\
\midrule
Novel algorithmic equivalences & No & Human proof \\
Concurrent/distributed equivalences & No & Human proof \\
Amortized complexity arguments & No & Human proof \\
\bottomrule
\end{longtable}

\begin{remark}
The ``Skeleton'' entries mean the tactic can set up the proof structure
(induction base/step, simulation init/step/output, etc.) but the
user must fill in the key insights---the invariant, the simulation
relation, or the abstraction function.  This is inherent: these
insights are the creative content of the proof.
\end{remark}


\section{Semi-Automated Proof Construction}

The typical workflow for proving a non-trivial algorithmic equivalence:

\begin{enumerate}
\item \textbf{Choose strategy}: based on the structure of the two
  programs, select simulation, abstraction refinement, induction, etc.

\item \textbf{Identify the key insight}: what is the simulation
  relation?  What is the abstract spec?  What is the loop invariant?

\item \textbf{Use a strategy tactic} to set up the proof skeleton:
  \texttt{by\_simulation(relation)}, \texttt{by\_spec(spec\_name)}, etc.

\item \textbf{Fill in sub-proofs}: for each subgoal, use
  $\mathsf{auto}$, $\mathsf{omega}$, $\mathsf{ring}$, or write
  explicit proof terms.

\item \textbf{Verify}: run the checker on the complete proof.  Fix
  any rejected sub-proofs.

\item \textbf{Serialize}: save the verified proof as a
  \texttt{.proof.json} file for permanent record.
\end{enumerate}

This workflow mirrors the practice in interactive proof assistants
(Lean, Coq, Isabelle) but is tailored to the domain of Python
program equivalence.


\section{LLM-Assisted Proof Construction}

An exciting application: LLMs can assist in proof construction.
Given two programs and a goal, an LLM can:

\begin{enumerate}
\item Suggest a proof strategy (simulation, abstraction refinement, etc.)
\item Identify the key relation or specification
\item Generate a proof term
\item The CCTT checker verifies the proof mechanically
\end{enumerate}

The LLM output is \emph{untrusted} until the checker verifies it.
This is the anti-hallucination principle applied to proofs:
the LLM's ``proof'' might be wrong, but the checker will catch
any errors.  The LLM is an oracle for proof \emph{search}, not
proof \emph{verification}.

\begin{remark}
This is a fundamental asymmetry: proof search is hard (creative),
but proof verification is easy (mechanical).  LLMs are good at
the creative part; the checker handles the mechanical part.
\end{remark}
